\begin{abstract}
A gyroscopic virtual mechanical system couples an undamped host oscillator to a
software oscillator through a skew-symmetric velocity interconnection.
For a coprime modal ratio $(\tau,\sigma)$ the closed loop evolves on closed
Lissajous orbits along which host energy can be transferred completely to the
controller.
This paper gives a complete account of the design, the failure, and the
restoration of this mechanism.
Exact modal--polar coordinates reduce the geometry to three quantities: total
energy, a modal amplitude ratio, and a coprime relative phase
$\psi=\sigma\theta_+-\tau\theta_-$.
In these coordinates, complete host absorption occurs if and only if the
amplitude ratio lies in an explicit band and the relative phase belongs to an
explicit finite set.
For equal amplitudes this set is $\psi\equiv(\tau+\sigma)(2j+1)\pi/2$.
The known condition $\tau-\sigma\equiv 2\pmod 4$ is exactly the case where the
impulse response satisfies it, and every resonant pair is absorbing at the
right phase.
We then show that the realised loop is fragile: for every positive loop delay the
interconnection is exponentially unstable, with a closed-form count of unstable
characteristic roots, the fast mode being destabilised and the slow mode damped;
a zero-order hold acts, to first order, as an additional loop delay of one
sampling period.
Retuning the virtual stiffness, as proposed earlier, can only compensate a
relative increase of the coupling of at most $n^2/8$, whereas a trim of the
common coupling has $O(1)$ authority.
A three-channel controller regulates energy, balances the modal amplitudes with
integral action, and locks the relative phase through the common coupling.
Its first-order averaged dynamics are triangular and input-to-state stable, and
semiglobal practical input-to-state stability of the sampled closed loop follows
by averaging.
Simulations with delay, sampling, gain errors, noise and backlash recover closed
orbits and per-cycle host depletion to $10^{-5}$ of the initial energy.
\end{abstract}

\begin{keyword}
Energy-based control; gyroscopic interconnection; vibration absorption;
time delay; sampled-data systems; averaging; input-to-state stability;
phase locking.
\end{keyword}
